% !TeX root = ../../Book4.tex
\chapter*{后记}
\addcontentsline{toc}{chapter}{后记}
\markboth{后记}{后记}

取整数 \(n\ge2\)，再对序言中的短正合列施行 \(\operatorname{Hom}_{\mathbb Z}(-,\mathbb Z)\)。末端留下的 \(\mathbb Z/n\mathbb Z\)，现在已经有了几种相互联系的解释。它是 Hom 消解复形的一次上同调，也是短正合扩张的等价类所组成的群，还可以写成导出范畴中的态射群
\[
\operatorname{Ext}_{\mathbb Z}^1(\mathbb Z/n\mathbb Z,\mathbb Z)
\simeq
\operatorname{Hom}_{D(\mathbf{Ab})}(\mathbb Z/n\mathbb Z,\mathbb Z[1]).
\]
这些解释分别回答计算、构造和比较的问题。一个由核、像得出的商群，逐渐获得了可以相加的扩张代表，也获得了与复形态射相容的表达。

\ProseParagraphBreak
消解把这个计算推广到更多对象。比较定理所做的工作尤其值得回顾：给出两套消解之后，先提升原对象间的映射，再证明不同提升只差一个同伦。于是计算结果既不依赖所选模型，又能随原来的映射自然变化。到了导出范畴中，拟同构被倒转，这些模型之间的比较成为真正的同构；消解仍然有用，它提供了可以动手计算态射的代表。

\ProseParagraphBreak
谱序列又把同一个计算分成若干层。四项双复形中，第一页留下两个非零项，第一次微分却为零；修正代表元以后，第二次微分把它们一起消去，与总复形的零调性相合。页稳定以后，还有另一件事要做：恢复滤过对象。\(\mathbb Z/4\mathbb Z\) 与 \((\mathbb Z/2\mathbb Z)^2\) 可以有相同的伴随分次，却保留不同的扩张。这两个例子让逐页计算与恢复对象的区别始终看得见。

\ProseParagraphBreak
导出范畴保留了不同次数之间的态射，而 \(t\)-结构决定怎样从中取出一颗心。标准心让我们重新认出原来的交换范畴；扭对倾斜则产生新的正合结构。\chapref{b4:ch:19}的两顶点例子把这种变化算得很具体：设 \(k\) 为域，\(\mathcal A\) 为箭图 \(1\to2\) 的有限维表示范畴，取简单对象 \(S_1=(k\to0)\)、\(S_2=(0\to k)\)。\cref{b4:prop:heart-higher-ext-mismatch}中的倾斜心 \(\mathcal B\) 满足
\[
\operatorname{Ext}_{\mathcal B}^2(S_1[1],S_2)=0,
\qquad
\operatorname{Hom}_{D^b(\mathcal A)}(S_1[1],S_2[2])\simeq k,
\]
这一差别说明，心的包含未必能延拓为全忠实的导出实现。判断这样的实现能否存在，需要比较心内的高阶 Ext 与环境中的移位态射。

\ProseParagraphBreak
最后的 DG 语言把态射本身组织成复形。在两个复形之间的 Hom 复形中，零次余圈给出复形映射，零次余边给出零伦映射；取零次上同调便回到熟悉的同伦类。保留整个 Hom 复形，还能继续追踪高次态射和产生同伦的资料。三顶点倾斜对象 \(T=P_1\oplus P_2\oplus S_2\) 给出的导出等价则有一个具体变化：\(\Phi=\mathbf R\operatorname{Hom}_R(T,-)\) 将 \(R=k(1\to2\to3)\) 的投射模 \(P_3\) 送到 \(B=k(1\to2\leftarrow3)\) 的简单模移位 \(S'_3[-1]\)。这个像只有一次上同调非零，所以等价没有保持两个模范畴的标准心；允许复形以后得到的等价，确实保存着不同于普通 Morita 等价的信息。

\ProseParagraphBreak
本卷中的定义很多。重读时，可以带着一个固定例子，把它送进 Hom、张量积或不变量函子，追问原来的正合关系在哪里改变了；再用消解、双复形或导出模型把改变算出来。这样，同一个问题会在不同语言中逐渐清楚。写这卷书所希望留下的，正是这种来回核对的习惯：抽象构造要能回到计算，计算中的现象也要能够提出新的数学问题。
